%!TEX root = hems_proj.tex
%%%%%%%%%%%%%%%%%%%%%%%%%%%%%%%%%%%%%%%%%%%%%%%%%%%%%%%%%%%%
\chapter{Fontosabb programkódok listája}\label{ch:progik}

\section{HEMS kódbázis}
A projekt forráskódja moduláris szerkezetű, és a felelősségek egyértelműen elkülönülnek az egyes fájlok között. A szcenáriók futtatását, az eredmények CSV/JSON formátumú mentését és az ábrák előállítását a \texttt{main.py} végzi, az optimalizáló eljárásokat, vagyis a GA, a PSO, a GWO és a hibrid GA--PSO implementációját az \texttt{algorithms.py} tartalmazza. A feladat fizikai modelljét, azaz a SmartHomeEnvironment-et, a célfüggvényt és a korlátkezelést a \texttt{hems\_problem.py} fogja össze, míg a szezononkénti fogyasztási, PV- és árprofilok generálásáért a \texttt{data\_generator.py} felel. Ez a tagolás teszi lehetővé, hogy a modell egy-egy eleme önállóan is érthető és módosítható legyen.

\section{Mintakódrészlet}
Az alábbi részlet a hibrid algoritmus fő lépéseit illusztrálja.
\lstinputlisting[language=Python, firstline=62, lastline=116]{../progfiles/algorithms.py}

\section{Modulszintű felelősségek}
A modulok felelősségi köre élesen elválik egymástól, ami a rendszer karbantarthatóságát szolgálja. A \texttt{hems\_problem.py} a célfüggvényt, a korlátkezelést, az energiamérleget és a töltöttségi (SOC) logikát tartalmazza, az \texttt{algorithms.py} a GA, a PSO, a GWO és a hibrid megoldó osztályait. A bemeneti oldalon a \texttt{data\_generator.py} állítja elő a szintetikus terhelés- és PV-profilokat, és látja el a CSV-k betöltését és validálását, az árakat pedig a \texttt{price\_loader.py} dolgozza fel az ENTSO-E adatok parse-olásával, az órás aggregációval és a napi ármap felépítésével. Mindezt a \texttt{benchmark\_real\_vs\_synthetic.py} hangolja össze, amely a profilépítéstől a kiértékelésen és az aggregáción át a riport és az ábrák mentéséig vezeti a folyamatot. Ez a szétválasztás csökkenti a modulok közötti csatoltságot, így egy-egy részegység, például az árbetöltés vagy a célfüggvény, a teljes pipeline újratervezése nélkül is továbbfejleszthető.

\section{Adat- és eredmény-artefaktok}
A benchmark kimenetei strukturáltan, több, egymásra épülő szinten jelennek meg, ami a teljes elemzés visszakövethetőségét biztosítja. A legalsó szinten a nyers, futásszintű táblázat áll, a \texttt{results\_realpv\_vs\_synth\_raw.csv} fájlban, amely minden egyes futás költségét és futási idejét rögzíti. Erre épül az aggregált összegzés, amely a summary, a gaps és a winners táblákban foglalja össze az algoritmusonkénti és adathalmazonkénti mutatókat, majd a szöveges markdown riport és a napi, illetve összegző ábrák zárják a sort. Mivel minden végső ábra és összegző érték visszavezethető a nyers táblázatokra, a dolgozatban szereplő állítások bármikor auditálhatók, és nem maradnak ellenőrizhetetlen, egyszeri benyomások.

\section{Reprodukálhatósági szempontok}
A kódban alkalmazott determinisztikus seedelés miatt ugyanazokkal a bemeneti fájlokkal a futások bármikor megismételhetők, a projektszerkezet pedig külön mappákban választja el a forráskódot, a riportot és az ábrákat, ami támogatja a verziózhatóságot és a visszaellenőrizhetőséget. Ennek a fegyelmezett szerkezetnek kutatási szempontból is jelentősége van, mert egyszerre garantálja, hogy az eredmények ne egyszeri, szerencsés futások legyenek, hogy a paraméterváltoztatások hatása futások között összehasonlítható maradjon, és hogy a későbbi bővítések, például a degradációs költség bevezetése vagy több eszköz párhuzamos kezelése, kontrollált körülmények között legyenek mérhetők.
